\documentclass[11pt,a4paper]{article}
\usepackage[T1]{fontenc}
\usepackage{lmodern}
\usepackage[margin=27mm]{geometry}
\usepackage{amsmath,amssymb,amsthm,booktabs,microtype}
\usepackage[hidelinks]{hyperref}
\newtheorem{proposition}{Proposition}
\newtheorem{lemma}[proposition]{Lemma}
\newtheorem{theorem}[proposition]{Theorem}
\theoremstyle{remark}
\newtheorem{remark}[proposition]{Remark}
\newcommand{\Ostar}{O^{*}}
\newcommand{\E}{\mathbb E}
\newcommand{\Pcal}{\mathcal P}
\newcommand{\Hbin}{H}
\emergencystretch=2em
\clubpenalty=10000
\widowpenalty=10000
\title{Exploring the Branching Frontier for Max 2-SAT\\
\large Full clause phases, joint completion bounds, and coverage}
\author{}
\date{7 October 2026}
\begin{document}
\maketitle

\begin{abstract}
We develop two concrete variations on the branching program of the October
5--6 notes. First, resolving a pruning clause completely gives an explicit
measure based on its exact completion count. Its analysis walk is uniform on
the permitted assignments, giving a simple concentration lemma. This changes
the algorithm and does not prove the earlier missing lemma for clauses that
are released. Second, counting the assignments satisfying a weighted sum of
active pruning clauses gives an overlap-aware, polynomial-time certificate
when the integer weights are polynomially bounded. Inherited certificates
give a branching inequality and an explicit variable-selection score. We
prove a conditional theorem identifying the cumulative progress still needed
for a fixed base below two. We also formalize when neighborhood coverage
helps, why it need not yield small separators, and how exact interaction
classes offer a different structural exit. No density-uniform running-time
theorem is claimed.
\end{abstract}

\section{Target and status}

The target is a deterministic polynomial-space algorithm with time
$\Ostar((2-\delta)^n)$ for a constant $\delta>0$ independent of a growing
density bound $\rho(n)$. Here $\rho=m_2/n$ counts binary clauses per variable;
the average occurrence degree is $d=2\rho$. An exponent saving that tends to
zero with $d$ does not meet this target, even if it is a useful improvement.
The first milestone could be all instances with $\rho(n)=\log\log n$;
$\rho=O(n)$ covers the full simple unweighted setting.

The notes already identify three difficulties: obtaining useful gaps,
overlap between pruning supports, and detecting contradictions before the
search reaches almost all variables. Our contributions below are mathematical
derivations within that program, with small exact checks. They have not been
independently refereed, and no novelty claim relative to all literature is made.

The polynomial-space baselines we checked have degree-dependent savings
\cite{Gol11,Has16}; these support the distinction above. Hastings also gives
instances with very many strict local optima. Consequently, a bound based only
on enumerating all local optima cannot provide a uniform constant saving.
Our web check did not locate a newer theorem settling the target; it is not
an exhaustive literature survey.

We use the deficit notation of the earlier notes. In a residual formula,
$p=\min(p_x,p_{\neg x})$, and $D$ is the free binary degree. For $B=1$,
branching against the literal with smaller deficit creates the valid
constraint $\langle L_x\le p-1\rangle$ on $D$ distinct variables. The preferred
child receives no new constraint. The soundness arguments in the earlier
notes apply to the modified strategies considered here.

\section{Resolve a pruning clause as one phase}

\subsection{An explicit branching rule and measure}

Define the exact cardinality volume
\[
 V(r,k)=\sum_{j=0}^{\min\{r,k\}}\binom rj,
 \qquad V(r,k)=0\quad\hbox{if }k<0.
\]
Consider this variant of the rule. At a phase boundary choose an undominated
$x$. Explore its preferred child recursively. In the other child enumerate
all assignments of its $D$ neighbors satisfying the new clause, then recurse
on the remaining variables. During this enumeration postpone other
domination fixings and additional pruning clauses. Earlier constraints may
be retained, but for the clean phase-distribution argument do not condition
the enumeration on them. They can be checked at the phase boundary.
This may do more work, but it preserves exactness. Combinations enumerate the
allowed assignments in polynomial space with polynomial overhead per output;
there is no need to scan all $2^D$ assignments.

\begin{proposition}[Full-phase branching]
For $k=p-1$, the phase has one child with a decrease of one free variable,
and $V(D,k)$ children with a decrease of $D+1$. It is admissible at base
$1<\lambda<2$ whenever
\begin{equation}
 \lambda^{-1}+V(D,k)\lambda^{-(D+1)}\le1,
 \quad\text{equivalently}\quad
 V(D,k)\le(\lambda-1)\lambda^D.                 \label{eq:phase}
\end{equation}
If every phase satisfies this condition, the search runs in
$\Ostar(\lambda^n)$ time and polynomial space.
\end{proposition}
\begin{proof}
These are exactly the enumerated alternatives. Induct on the number of free
variables using \eqref{eq:phase}. The work enumerating phase alternatives
is polynomial per child, and there are at most $n$ levels of recursive
phases on a path. Retained constraints can only remove alternatives.
\end{proof}

For an equivalent binary-step analysis put
\begin{equation}
 e_\lambda(r,k)=r-\log_\lambda V(r,k),\qquad
 \mu=u+\sum_C\log_\lambda V(r_C,k_C)
       =n_{\rm free}-\sum_C e_\lambda(r_C,k_C),                 \label{eq:signed}
\end{equation}
where $u$ counts uncovered free variables and the counted supports are
disjoint. A single active clause is sufficient for the phase algorithm.
Weights $e_\lambda$ are permitted to be negative. This is intentional:
a weak residual constraint carries the cost of completing its support.
The measure is still nonnegative because each live clause has $V\ge1$.

\begin{lemma}[Binary steps and forced assignments]
Resolving a clause variable satisfies the branching inequality with equality
for the measure \eqref{eq:signed}. A forced assignment never increases
$\mu$. These statements do not require $e_\lambda\ge0$.
\end{lemma}
\begin{proof}
Use $V(r,k)=V(r-1,k)+V(r-1,k-1)$. If the literal is set to $b$, the measure
drop is $\log_\lambda(V(r,k)/V(r-1,k-b))$. Summing $\lambda^{-\text{drop}}$
over the nonempty children gives one. Restriction cannot increase $V$.
For an uncovered variable the drop is one. At $k=0$ there is one completion,
so propagating the remaining literals also respects the measure.
\end{proof}

The same inequality shows that a transition cannot decrease $e_\lambda$
by more than one. At a vacuous state $e_\lambda=r(1-\log_\lambda2)<0$.
Keeping the clause until its support is completed is permitted. Releasing
it would reduce the measure, but would change the distribution used next.

\subsection{A concentration lemma without an unproved walk assumption}

\begin{proposition}[Distribution inside an isolated full phase]
Choose each child during clause resolution with probability equal to its
remaining completion count divided by that of its parent. The resulting
assignment is uniform among the $V(D,k)$ permitted assignments. If $J$ is
its number of true pruning literals and $0\le k<D/2$, then for $1\le t\le k$,
\begin{equation}
 \Pr[J\le k-t]\le
 \left(\frac{k}{D-k+1}\right)^t.                 \label{eq:tail}
\end{equation}
\end{proposition}
\begin{proof}
The child probabilities telescope to $1/V(D,k)$ at every completion.
For $q=k/(D-k+1)$ and $j\le k$,
$\binom D{j-1}/\binom Dj=j/(D-j+1)\le q$. Thus
$\binom D{i}\le q^t\binom D{i+t}$ for $i\le k-t$.
Sum and divide by $V(D,k)$.
\end{proof}

For fixed $k/D\le\beta<1/2$, deviations of $\eta D$ below the allowance
have probability at most $\exp(-\Omega_\beta(\eta D))$. This proves a
phase-level statement at least as strong as the sought concentration scale,
\emph{for this modified full-resolution rule}. It does not establish that
statement for the October 6 walk, whose zero-weight states are released.
It also does not establish the availability of subsequent good phases.

The mean satisfies
\[
 \E[J]\ge k-\frac{q}{1-q}.
\]
By symmetry each support literal is true with probability
$\E[J]/D=k/D-O_\beta(1/D)$ when $k/D\le\beta<1/2$.
For another variable, its helpful values need not agree with all of these
literals, and it may share only part of the support. The following bound
handles both issues for a single complete phase.

\begin{lemma}[Helpful updates on any fixed overlap]
Fix $s\ge1$ variables of the clause support and designate one helpful value
for each, with arbitrary signs. Let $Y$ be the number taking that value in
a uniform permitted completion. If $0<\eta<k/D<1/2$ and
$q=k/(D-k+1)$, then
\begin{equation}
 \Pr\!\left[Y<\left(\frac{k}{D}-\eta\right)s\right]
 \le q^{\lceil\eta D/2\rceil}+2e^{-\eta^2s/8}.       \label{eq:overlap}
\end{equation}
\end{lemma}
\begin{proof}
By \eqref{eq:tail}, except with probability at most the first term,
$J/D\ge k/D-\eta/2$. Condition on such a value $J=j$ and write $p=j/D$.
The true positions form a uniform $j$-subset. Split the designated positions
into $a$ where true is helpful and $b=s-a$ where false is helpful.
Their helpful counts have means $pa$ and $(1-p)b$, whose sum is at least
$ps$, since $p<1/2$. Each count is a sample without replacement from a
zero-one population. Hoeffding's bound \cite{Hoe63} applies to each.
For their sum to be below $(k/D-\eta)s$, one of the two counts must be at
least $\eta s/4$ below its mean. The two lower-tail bounds are at most
$e^{-\eta^2s^2/(8a)}$ and $e^{-\eta^2s^2/(8b)}$, omitting an empty group.
Each is at most $e^{-\eta^2s/8}$. A union bound proves the claim; no
independence between the two counts is assumed.
\end{proof}

The designation is fixed at the beginning of the phase, for example toward
a chosen deficit of a variable that remains outside its support. This is
conditional on the phase-start history and imposes no alignment assumption.
It is not a theorem about many adaptive phases: their ratios $k/D$ may be
too low, intersections may be small, and the deficit orientation may change.

\subsection{Entropy scale and a remaining availability problem}

For $k\le\beta D$ with $0<\beta<1/2$,
$V(D,k)\le2^{H(\beta)D}$. Therefore, for any
$2^{H(\beta)}<\lambda<2$, every such phase with
\begin{equation}
 D\ge \frac{\ln(1/(\lambda-1))}
                 {\ln\lambda-H(\beta)\ln2}                 \label{eq:threshold}
\end{equation}
is feasible. Small degrees can be checked directly or handled separately.
Do not infer uniform finite-size feasibility at the exact entropy base
from an asymptotic estimate alone.

For $k=D/4$, the numerical roots of \eqref{eq:phase} are:
\begin{center}
\begin{tabular}{rrrrr}
\toprule
$D$ & 16 & 32 & 64 & 128\\
\midrule
full-phase base & 1.67236 & 1.69519 & 1.71536 & 1.73019\\
\bottomrule
\end{tabular}
\end{center}
They approach $2^{H(1/4)}\simeq1.75477$. These are local branching numbers,
not global bounds without an availability argument. The earlier optimized
release rule can have smaller finite-size branching numbers.

Allow plain steps and use the previous averaging construction. At a feasible
phase give the preferred child probability $1/\lambda$ and spread the
remaining probability uniformly across its $V(D,k)$ alternatives. Each
alternative then has probability at least $\lambda^{-(D+1)}$ by
\eqref{eq:phase}. The proof of the earlier averaging lemma gives
\[
 L\le\lambda^n\E[(2/\lambda)^S]
\]
for the number $S$ of plain steps on this phase walk. The exponentially
weighted cost of plain steps remains the actual obstacle. A small empirical
mean for $S$ is insufficient.

A useful next experiment is to choose the largest feasible ratio below a
fixed threshold, rather than the smallest ratio. A very strong clause can
force values unfavorable to its neighbors; a ratio near the threshold
passes on a less biased distribution. This is a motivated strategy, not
a proved universal improvement.

\section{A joint measure from exact weighted-sum counts}

\subsection{The certificate}

At a search state let the active pruning inequalities be
$L_i(x)\le k_i$, $i=1,\ldots,m$, over $r$ free variables. For nonnegative
integer weights $a_i$ define
\begin{equation}
 N_a=\#\{x\in\{0,1\}^r:\ \textstyle\sum_i a_i L_i(x)
                          \le\sum_i a_i k_i\}.             \label{eq:Na}
\end{equation}
Every completion satisfying all active clauses is counted by $N_a$.
Unlike an additive clause weight, $N_a$ explicitly incorporates overlapping
supports, cancellation and redundant clauses.

For free variable $w$, let $A_w^b$ be its contribution to
$\sum_i a_iL_i$ when it takes value $b$. Set
\[
 K=\sum_i a_i k_i-\sum_w\min(A_w^0,A_w^1),\qquad
 c_w=|A_w^1-A_w^0|.
\]
After orienting each variable so that its cheaper value is zero,
\begin{equation}
 N_a=\sum_{j\le K}[z^j]\prod_w(1+z^{c_w}).                  \label{eq:dp}
\end{equation}
This is a one-dimensional integer dynamic program. It takes
$O(r(1+\sum_w c_w))$ arithmetic operations, with counts having at most
$r+1$ bits. If clause multiplicities and the weights $a_i$ are bounded by
polynomials in the input size, both time and space are polynomial. Merely
bounding the \emph{bit length} of arbitrary huge weights is not enough.

The unweighted sum cut is exactly the test $N_{\boldsymbol 1}=0$.
The certificate gives useful information before that cut becomes possible.
For a single cardinality clause it gives its exact binomial volume, including
the roughly one-half factor for a nearly balanced clause that a Chernoff
bound can miss.

\subsection{Inheritance makes the branching inequality rigorous}

\begin{theorem}[Joint completion certificate]
Store a weight vector $a$ and $U=N_a$ at every live node. Cut if $U=0$.
For each child, evaluate the inherited vector after substitution and any
polynomial-size list of additional candidate vectors, and store the smallest
count. Extend the inherited vector by zero on newly created pruning clauses.
Then, after branching and any further sound fixings,
\begin{equation}
 U_0+U_1\le U.                                             \label{eq:super}
\end{equation}
Also $U\ge1$ at every kept node and $U=2^n$ at an unconstrained root.
\end{theorem}
\begin{proof}
For a fixed vector, assignments counted at the parent partition according
to the value of the branching variable:
$N_a=N_{a,0}+N_{a,1}$. Further fixings restrict these sets. A new constraint
has inherited weight zero and therefore does not invalidate the old
certificate. Minimization over a list containing the inherited vector
cannot increase either child count. These facts prove \eqref{eq:super}.
The counts are nonnegative integers and upper-bound all jointly feasible
completions, so a zero count justifies a cut.
\end{proof}

The value $U$ here is part of the augmented search state, together with its
witness. We do not claim that globally minimizing \eqref{eq:Na} over all
bounded integer vectors is polynomial-time. Polynomial candidate lists and
bounded rounds of coordinate improvement suffice for the theorem.
If every state were optimized independently by a heuristic without retaining
the inherited witness, \eqref{eq:super} would not be guaranteed.

Useful candidates are the zero vector, each single clause, all clauses with
weight one, promising clause pairs, and bounded coordinate modifications.
A linear program or a smooth Chernoff optimization can propose weights;
the final integer count still certifies the result exactly.

\subsection{A concrete strategy and the missing global theorem}

The measure $\log_2 U$ has a valid base-two branching inequality. It does
not yet give a base below two. For candidate variable $x$, form its two
children, including its new pruning clause, and calculate
\begin{equation}
 I_x=\log_2 U-\log_2(U_{x=0}+U_{x=1})\ge0.                \label{eq:info}
\end{equation}
Choose the variable maximizing this certified decrease. If the denominator
is zero, both children can be discarded. Polynomially many candidates and
certificate evaluations preserve polynomial time per visited node and
polynomial space. This is an explicit branching strategy optimizing a
proved certificate; its worst-case advantage is still unproved.

It is essential to count work spent forming all candidate children. A
smaller selected search tree need not give a faster implementation at small
sizes. Fixed polynomial lookahead overhead is harmless for an eventual
exponential-base theorem, but matters in experiments.

\begin{theorem}[A sufficient cumulative-progress condition]
Use the certificate search above, with polynomial work per visited node.
Normalize child probabilities to $U_i/(U_0+U_1)$ whenever the sum is nonzero.
Absorb successful terminal nodes with $U=1$, and regard a node with no kept
children as death. Let $G_t$ be the sum of \eqref{eq:info} along this walk
through branching depth $t$, with $G_t=+\infty$ after death.

Suppose, for constants $\alpha,\gamma,\xi>0$ and
$t_0=\lfloor(1-\alpha)n\rfloor$,
\[
 \Pr[G_{t_0}<\gamma n]\le2^{-\xi n}.
\]
Then the search runs in
\[
 \Ostar\!\left(2^{(1-\min\{\alpha,\gamma,\xi\})n}\right).
\]
\end{theorem}
\begin{proof}
Write $R_v=(U_{v0}+U_{v1})/U_v=2^{-I_v}$. For a frontier node $w$, the
probability of reaching it under normalized transitions times
$\prod_{v\prec w}R_v$ is $U_w/2^n$. Hence the total certificate mass at
depth $t$ is $2^n\E[2^{-G_t}]$.
At $t_0$ this is at most
$2^n(2^{-\gamma n}+2^{-\xi n})$.
Subsequent total mass cannot increase, and each kept node has integer mass
at least one. There are at most $n$ branching levels. The earlier levels
have at most $2^{t_0+1}$ nodes. Discarded children and polynomial lookahead
add only polynomial overhead to these counts.
\end{proof}

This is a precise replacement for ``encounter many variables with gaps''.
The required quantity is linear cumulative \emph{joint} progress, early
enough in the tree, with exponentially controlled bad paths under the stated
walk. Repeatedly creating almost redundant clauses need not contribute to it.
So far no condition depending only on $\rho(n)=\omega(1)$ proves this premise.

\subsection{A real limitation of weighted sums}

Suppose the relaxation $0\le x_w\le1$, $L_i(x)\le k_i$ is feasible.
For any nonnegative weights $a$, its feasible fractional point gives
\[
 \min_{x\in\{0,1\}^r}\sum_i a_i L_i(x)
 =\min_{x\in[0,1]^r}\sum_i a_i L_i(x)
 \le\sum_i a_i k_i.
\]
The equality holds because the objective is linear on a box. Therefore
$N_a\ge1$ for every weighted sum, even if the Boolean conjunction is
infeasible. For example, $x+y=y+z=z+x=1$ has the fractional solution
$(1/2,1/2,1/2)$ but no Boolean solution. Encode each equality by its two
cardinality inequalities. No nonnegative weighted-sum cut detects the
contradiction at the root.

Conversely, separating the image of the box from the lower orthant shows
that an infeasible linear relaxation admits a nonnegative combination whose
minimum exceeds its bound. Such multipliers need not lie in our small
candidate list or have small numerical values. Thus weighted sums need
either branching or stronger integer reasoning in the remaining cases.

\section{What coverage actually certifies}

Let $C$ be the union of selected active supports and let
$S\subseteq V\setminus C$ be a set of candidate variables.

\begin{proposition}[Exact and approximate coverage]
If $N(v)\subseteq C$ for every $v\in S$, then $S$ is independent and the
residual Max 2-SAT instance can be solved in time $\Ostar(2^{n-|S|})$ and
polynomial space. More generally, if every $v\in S$ has at most $b$ neighbors
outside $C$, then a polynomial-time construction gives an independent set
of size at least $|S|/(b+1)$, and time
$\Ostar(2^{n-|S|/(b+1)})$ follows.
\end{proposition}
\begin{proof}
In the first case no edge has both endpoints in $S$. Assign its complement
and optimize every variable of $S$ independently. In the second case
$G[S]$ has maximum degree at most $b$; a greedy independent-set construction
removing a chosen vertex and its neighbors proves the claim.
\end{proof}

Here $n$ is the number of variables in the residual subproblem. The statement
does not require a small separator. For instance, $K_{t,t}$ has exact
coverage of one side by the other, a linear independent set, and no
sublinear balanced separator. Conversely, if $C$ contains nearly every
variable, coverage by $C$ says little. Also, $b=\varepsilon d$ is not enough
for a constant saving when $d$ grows: the independent-set guarantee can
still be sublinear.

A useful coverage dichotomy therefore must quantify both the number of
\emph{uncovered centers} and the absolute amount of their remaining internal
interaction. Merely counting covered neighbors, or measuring a covered
fraction, misses these quantities.

If a candidate cut $C$ really does leave components $K_j$ with
$\max_j|K_j|=q$, independent optimization after enumerating $C$ gives time
$\Ostar(2^{|C|+q})$. Inherited pruning clauses can couple these components.
One safe terminal action is to discard those auxiliary constraints and solve
the residual \emph{soft} Max 2-SAT objective exactly. Every returned assignment
is still an assignment of the original formula, and the enlarged subtree
contains every originally retained completion. This is valid for this
optimization search, though it forfeits the discarded pruning information.

\section{An alternative structural exit: interaction classes}

Small separators are not the only useful outcome of aligned neighborhoods.
Write a residual objective in Boolean quadratic form
\[
 f(x)=c+\sum_i h_i x_i+\sum_{i<j}J_{ij}x_i x_j.
\]
Suppose the remaining variables can be partitioned into $t$ classes such
that $J_{ij}$ depends only on the classes of $i,j$, including a uniform
coefficient for distinct variables in the same class. The $h_i$ can differ.

\begin{proposition}[Exact interaction classes]
Such an instance is solvable in polynomial space and time
$\Ostar(\prod_{a=1}^t(|C_a|+1))\le\Ostar((n+1)^t)$.
If the property holds after deleting an exception set $E$, time
$\Ostar(2^{|E|}(n+1)^t)$ suffices.
\end{proposition}
\begin{proof}
Enumerate the number $k_a$ of ones in each class. All quadratic contributions
are determined: they are $J_{aa}\binom{k_a}{2}$ within a class and
$J_{ab}k_a k_b$ across classes. The best linear contribution selects the
$k_a$ largest $h_i$ in each class. Enumerating $E$ changes only fields on
the remaining vertices and a constant.
\end{proof}

Thus $|E|\le(1-\delta)n$ and $t\log_2(n+1)=o(n)$ give a fixed-base
improvement. This extends the complete aligned example in the earlier notes
and can apply with very large separators. It requires \emph{exact signed
interaction classes}. Shared supports, large Gram correlations, approximate
low rank, and a small number of active rows do not alone prove this property.
Turning approximate alignment into a small exception set or a manageable
residual interaction graph is a separate open direction.

\section{Checks and what the experiments do not establish}

The reproducible script is \path{explore_branching_frontier.py}.
Its output is \path{branching_frontier_checks.json}. No existing notes or experiments
were modified. The script uses exact integer counts, except for reported
branching roots and numerical checks of logarithmic identities.

\begin{itemize}
\item The full-phase search matches brute-force optimum on 80 random
$B=1$ formulas with $3\le n\le10$, exercising 103 phase nodes.
\item There are 41,650 checks of \eqref{eq:tail} and 3,599 checks of forced
measure transitions, including negative-weight states.
\item On 150 random cardinality systems with $3\le n\le9$, every tested
certificate upper-bounds the brute-force joint completion count. All 918
variable-splitting identities and inequalities pass, as do 150 tests adding
an extra constraint in one child.
\item The small candidate optimizer improves over the best single-clause
bound in 55 of these 150 systems, and over the all-ones weighted sum in 103.
These are certificate comparisons on synthetic systems, not Max 2-SAT
running-time estimates.
\item A contradictory pair on seven variables plus a weak clause on 20 new
variables has all-ones count $134{,}215{,}040$. Weights $(1,1,0)$ give zero.
The odd-cycle equality example remains positive for every tested vector
with weights in $\{0,1,2\}$; its actual joint count is zero.
\end{itemize}

We also implemented one-step candidate branching using \eqref{eq:info} and
compared it with static branching on six random $B=1$ formulas, edge
probability $0.7$. Both versions use the same completion certificates and
domination. For speed, their weight candidates are just zero, singletons,
all ones, and the inherited vector. All optimum values match brute force.

\begin{center}
\begin{tabular}{rrrrr}
\toprule
& \multicolumn{2}{c}{kept nodes visited}
& \multicolumn{2}{c}{candidate children evaluated}\\
$n$ & static & volume strategy & static & volume strategy\\
\midrule
12 & 14 & 6 & 22 & 56\\
12 & 16 & 7 & 28 & 86\\
12 & 4 & 2 & 4 & 22\\
16 & 91 & 32 & 162 & 694\\
16 & 40 & 19 & 72 & 412\\
16 & 107 & 51 & 190 & 866\\
\bottomrule
\end{tabular}
\end{center}

The selected trees are smaller, but the lookahead performs more work.
These tiny instances establish implementability and provide correctness
checks; they do not show an asymptotic improvement, nor even a wall-clock
speedup. They are also too small to test a growing-density claim.

\section{The next mathematical questions}

\paragraph{Availability under complete phases.}
Use the single-phase overlap bound \eqref{eq:overlap} to study drift across
adaptively chosen phases. Investigate largest-feasible-ratio branching and
bound the exponential moment of plain steps. The unresolved issue is to
control the chosen clause ratios and accumulation of local errors; full
phases do not provide this global drift theorem for free.

\paragraph{Joint progress before the last levels.}
Prove the cumulative-progress premise in the theorem for an interesting
class, initially pseudorandom signed instances. The critical quantitative
question is whether useful inequalities add linear total information before
branching depth $(1-\alpha)n$. Uniform assignments and the certificate walk
are different distributions; substituting one for the other is not valid.

\paragraph{From coverage failure to a certified structural exit.}
Track the size of the uncovered candidate set, its internal degree, the
components after a prospective cut, and signed interaction classes. A valid
dichotomy must either create genuinely new joint progress or provide a
quantified residual solver. High overlap alone cannot certify either side.

\paragraph{Beyond one aggregate inequality.}
The LP-feasible but integer-infeasible example motivates small groups of
aggregate inequalities, parity reasoning, or class decisions. A constant
number of weighted inequalities can be counted by a multidimensional dynamic
program when all coefficient sums are polynomial. Letting this dimension
grow raises a polynomial-space issue; the resulting table must not be
silently charged as polynomial storage.

\paragraph{The density milestone.}
For a theorem covering all formulas with a growing density bound, explicitly
include $B=2$ instances such as Max-Cut encodings. The weighted-sum certificate
extends to integer multiplicities directly; the exchangeable binomial phase
lemma is specifically $B=1$. A weighted phase has a computable generating
function, but no automatic exchangeability. The missing result is a
structural/probabilistic argument covering every instance at the selected
growing density. None of the checks above supplies it.

\begin{thebibliography}{9}
\bibitem{N3} Max 2-SAT and the domination rule: results and ideas so far.
Working note, 5 October 2026.
\bibitem{N4} Max 2-SAT with pruning clauses: pair decisions, sums, and a
measure. Working note, 6 October 2026.
\bibitem{Gol11} A. Golovnev. New upper bounds for MAX-2-SAT and MAX-2-CSP
w.r.t. the average variable degree. IPEC 2011.
\url{https://golovnev.org/papers/max2sat.pdf}.
\bibitem{Has16} M. B. Hastings. Local maxima and improved exact algorithm
for MAX-2-SAT. 2016. \url{https://arxiv.org/abs/1610.07100}.
\bibitem{Hoe63} W. Hoeffding. Probability inequalities for sums of bounded
random variables. \emph{Journal of the American Statistical Association}
58(301):13--30, 1963. \url{https://doi.org/10.1080/01621459.1963.10500830}.
\end{thebibliography}
\end{document}
